%% Section 1: Introduction. Jobs: field stake -> bottleneck -> prior
%% attempts -> what this paper does -> contributions. (outline v2 sec.3.1)

Streaming speech services are increasingly built on \emph{revision-generative}
models: a Whisper-style decoder is conditioned on its own previously generated
text~\cite{radford2023whisper,machacek2023whisperstreaming}, and ASR-capable
large language models such as Qwen3-ASR re-prompt themselves with rollback
prefixes of their prior output~\cite{shi2026qwen3asr}. In these systems the
generated text is not merely displayed; it is written back into the state that
conditions future generation. Serving systems, however, are still built on the
opposite assumption: a generated output is terminal, and its only consumer is
the user~\cite{kwon2023vllm,yu2022orca,zheng2024sglang}.

Batching is the standard mechanism by which serving systems buy throughput:
independent sessions share one forward pass~\cite{kwon2023vllm,agrawal2024sarathi,zhong2024distserve,sun2024llumnix}.
But batched execution perturbs floating-point arithmetic relative to singleton
execution, so batched outputs can differ locally from their canonical
single-session counterparts. In terminal-output systems this difference is
display noise. In revision-generative systems it is worse: once a perturbed
output is written into feedback state, the divergence is inherited by every
future revision. Section~\ref{sec:model} makes this precise; our causal
experiment confirms the prediction: with the feedback edge present, 8 of 21
legal state pairs propagate a batch-created divergence, and with the edge
removed, none do; on the persistent pairs the measured revision amplification
factor (RAF) reaches 153.

Prior mechanisms do not cover this failure mode. Streaming-ASR confirmation
policies such as LocalAgreement commit stable display prefixes but never
constrain which outputs may become future-state
parents~\cite{machacek2023whisperstreaming}. Speculative decoding verifies
draft tokens against a target distribution, not state-write
authority~\cite{leviathan2023specdec,chen2023specsample}. Deterministic-inference
work repairs arithmetic reproducibility but not feedback
semantics~\cite{xu2022deterministic,antunes2024reprohpc}. Optimistic
concurrency control supplies the closest analogy~\cite{kung1981occ,bernstein1983mvcc},
yet classical transactions commit reads and writes of stored data, whereas
here the ``write'' is a model-generated projection whose legality must be
certified by the serving contract itself.

This paper takes one position: \emph{revision outputs are speculative writes
to future model state, not merely user-visible responses}, and
\[
\mathrm{Display}(\out_t)\;\not\Rightarrow\;\mathrm{Authorize}\bigl(\proj(\out_t;\state_t)\bigr).
\]
We introduce \textbf{Dual-Granularity Feedback Transactions} (DGFT): a fine
proposal path serves low-latency display, while only a coarse canonical
authority may install state projections, and only as a certificate-bearing
transaction. DGFT is formalized in Section~\ref{sec:dgft}, where
same-granularity certification appears as the degenerate $\granratio=1$ case,
three transaction invariants are proved, and an authority-cost decomposition
predicts when coarse authority outperforms same-granularity verification.
\textbf{RevTrace} (Sections~\ref{sec:system}--\ref{sec:impl}) implements the
transaction runtime with a stratified, contract-scoped certification layer.

Contributions:
\begin{enumerate}
  \item A problem characterization of feedback-state contamination in
    revision-generative serving, with a causal verification that the feedback
    edge is the propagation channel (Section~\ref{sec:model},
    Section~\ref{sec:eval}, H1).
  \item DGFT, a dual-granularity transaction mechanism with provable
    invariants---uncertified-ancestry exclusion, committed-state
    linearizability, and granularity-independent safety---and a cost analysis
    whose predictions we test (Section~\ref{sec:dgft}, H3).
  \item RevTrace, a runtime with contract-scoped stratified certification:
    randomized audits and anytime-valid bounds on the same-granularity path,
    FP32 preflight invariance checks, exhaustive transaction validation, and
    independent quality references (Section~\ref{sec:system}, H2, H5).
  \item A cross-family evaluation on two architecture families and five
    checkpoints (Qwen3-ASR 0.6B/1.7B; Whisper tiny/base/small.en) on two
    RTX 3090 GPUs: same-workload SG/Coarse/Hybrid ablation, production
    performance, per-family quality contracts, and boundary evidence
    (Section~\ref{sec:eval}, H1--H5).
\end{enumerate}
